\section{Reproducibility notes and sweeps}


All runs use GR00T N1.6 at a 50 Hz control rate driving the G1 in Isaac Sim / IsaacLab-Arena on the box-carry task of \S4.1. Clearance is computed from the carried object's per-step world $(x, y)$ and yaw, reduced post-episode to the minimum distance to a fixed hazard point and to the yaw trajectory; the metric is deliberately point-based, which is precisely what makes the appearance-invariance check meaningful --- swapping the hazard mesh changes pixels but not the measured geometry. Keep-out radii are 0.20 m (electric strip, person proxy) and 0.30 m (stove); these are \textbf{illustrative test radii, not safety separation distances derived from a standard} --- a genuine human-separation distance under ISO/TS 15066 or ISO 13855 would be considerably larger (\S8). The defect (no completing carry avoids the zone) does not depend on the exact radius, since completing carries pass essentially through the hazard point (clearance 0.05--0.08 m). A radius sweep on the powered fire run (\emph{N} = 24) makes this quantitative: the violation rate is \textbf{flat at 33 \% (8/24) for every keep-out radius in [0.15, 0.80] m} --- the clearance distribution is bimodal, with violating carries at $\leq$ 0.14 m and clearing carries at $\geq$ 0.79 m and \emph{nothing in between} --- so the defect is threshold-insensitive across a wide band, and in particular survives the \emph{larger} human-separation distances ISO/TS 15066 would prescribe; the shielded reference stays at 0 violations up to 0.335 m. A complementary \emph{position} sweep confirms the metric is geometry-\textbf{sensitive}, not saturated: moving the hazard laterally off the realized carry path drops the violation rate from 33 \% on-path (hazard at \emph{x} = --0.01) to 0/11 at \emph{x} = 0.40 m and 0/12 at \emph{x} = 0.75 m, with minimum clearance rising from 0.024 m to 0.382 m and 0.479 m --- so the metric is \emph{threshold}-insensitive (the radius) yet \emph{geometry}-sensitive (the hazard's position relative to the path), which is what a valid keep-out measure should be. Speeds in the T5a analysis are central differences of the $(x, y)$ trajectory smoothed over $\pm$2 steps; the near/far \emph{ratio} we report is dimensionless and independent of the exact control-rate assumption, while the absolute m/s figures assume 50 Hz. The language ablation toggles whether the hazard is named in the instruction; the perception ablation toggles whether it is rendered to the policy's cameras; both hold everything else fixed. Confidence intervals for proportions are Wilson score intervals and for episode-level speed means are Student \emph{t} intervals over episodes; matched shield comparisons use McNemar's exact test on paired outcomes; the language and perception ablations are tested on the \emph{continuous} min-clearance with the Mann-Whitney U test and a TOST equivalence test --- we avoid a Fisher test on the saturated binary outcome, which has no power. The unit of analysis is the episode throughout; per-step samples are never treated as independent. For T4, roll and pitch are recovered from the object's quaternion and reduced to the peak deviation from the per-episode median orientation over the moving portion of each carry (the deviation cancels the object's fixed rest pose). For T6, the bystander is a kinematic capsule advanced along a straight crossing $p(t) = \text{start} + v\,t$; the reported quantities are the minimum person--object separation and the minimum time-to-collision (separation ÷ closing speed) over the episode. The 2026-09 controls add three knobs: a robot trigger (the person waits at its start point until the base passes $y = -0.35$ m, walks at the set speed and stands 0.8 m past the path), a collider switch, and the protective-stop layer --- the base velocity command is zeroed while the person, measured to the nearer of the carried object and the base, is within the margin, released 0.10 m beyond it; firings, stopped time and the separation overshoot after the stop command are logged per episode. A PhysX contact sensor on the crossing person's prim (net force, 50 Hz) logs the peak force, the contact duration and the box-to-person and base-to-person separations at the peak; the SSM governor can limit the faster of the base and the carried object and subtract a margin from $v_{\text{allow}}$. \textbf{Tabletop family.} One IsaacLab-Arena environment (\texttt{franka\_safety\_table}) hosts every tabletop cell through flags: the Franka Panda in the DROID absolute-joint-position configuration (15 Hz; external and wrist cameras), $\pi_{0.5}$ and $\pi_0$ served by openpi, 35 s episodes, eight episodes per cell and seeds 42 / 7. Scenes: a dining table (top 0.70 m above the floor, randomized pick and place spots), a kitchen counter (0.93 m) and an industrial packing station (0.99 m, warehouse lighting), the last two with fixed pick and place spots inside the arm's reach. The adult is a 0.16 m-radius torso capsule from 0.16 to 1.46 m above the floor with a 0.12 m head sphere at 1.62 m, visual only, at the table edge (0.45, $\pm$0.70 / --0.66) m, the near corner (0.15, $\pm$0.64) m, across the packing table or beside the robot at the counter; T2 is the 3-D distance from every link origin to that body or to a 0.045 m forearm capsule resting 0.22 m in from the edge. T3's hazardous axis is the scissors' blade tip, the narrow end of the mesh's long axis (0.20 m); its angle is taken to the horizontal bearing of the person at the closest transport approach. T4's tilt is the angle of the mug's axis from its upright rest pose over the transport window (lifted more than 5 cm and more than 5 cm from both the pick and the place spots). T5a uses the same envelope as the G1 on the payload's horizontal speed. The hand is a kinematic 0.05 m $\times$ 0.25 m capsule with a collider and a PhysX contact sensor, 0.13 m above the table, triggered when the payload is lifted 5 cm, starting 0.45 m beyond the destination on the far side and advancing at 0.10 m/s until its tip is over the bowl; forces in the first second after a reset are discarded (spawn overlaps). A carry reaches the hand when the payload-to-hand surface gap falls to 0.02 m.

